%% Abstract (limit 200 words, also submitted on 18 Jan 2027).
LLM-driven algorithm design uses a language model as a black-box
variation operator: unlike for a classical operator, we do not know what
it knows, what it uses, or how to control it. We open this operator with mechanistic interpretability,
reading and editing the model's activations while it writes optimisation
heuristics in \llamea. \emph{What it knows:} whether the new algorithm
will crash is decodable while the code is written (AUROC 0.85, +0.08
over code text, output uncertainty and search metadata), also on new
problems, and sparse-autoencoder features represent metaheuristic
concepts beyond their keywords. \emph{What it listens to:} in three
models, the edits ignore the score in the feedback and follow the parent
code and error traces. \emph{How to control it:} a direction read from
the activations acts as an exploration dial that changes the measured
search behaviour and, on five unseen problems, the algorithm families
written. In the loop, steering towards exploitation narrows the design
space and costs up to 0.24 \aocc; a dose adapted to the success of the
offspring, like a step size, is the best arm on MA-BBOB. A pre-registered
test on BBOB 20-D shows that the dial acts on the design space, not on
the optimiser's behaviour.
